\documentclass[10pt,a4paper]{amsart}
\usepackage[margin=19mm]{geometry}
\usepackage{amsmath,amssymb,mathtools}
\usepackage[scr=rsfs,cal=euler]{mathalfa}
\usepackage{xfrac}
\usepackage[unicode,hidelinks,bookmarks=false]{hyperref}
\newcommand{\ie}{i.e.\ }
\newcommand{\cf}{cf.\ }
\newcommand{\resp}{respectively\ }
\newcommand{\Subsection}{\subsection}
\providecommand{\nobreakdash}{\nobreak\hbox{-}}
\setlength{\emergencystretch}{2em}
\multlinegap0pt
\usepackage{amssymb}
\usepackage[shortlabels]{enumitem}
\usepackage{xcolor}
\usepackage{tikz}
\usepackage{mathtools}
\newtheorem{proposition}{Proposition}
\def\NN{\mathbb{N}}
\newcommand\restr[2]{\ensuremath{\left.#1\right|_{#2}}}
\title{Information Geometry via the $Q$-Root Transform}
\author{Levin Maier}
\date{Source: arXiv:2603.20081v1 (CC BY 4.0)}
\begin{document}
\maketitle

\noindent\emph{Source and licence.} Information Geometry via the $Q$-Root Transform. Version arXiv:2603.20081v1 (CC BY 4.0), distributed by arXiv under the Creative Commons Attribution 4.0 International licence, http://creativecommons.org/licenses/by/4.0/, as recorded in the arXiv licence metadata for this exact version and shown on the abstract page https://arxiv.org/abs/2603.20081v1. Full licence notice: \texttt{licenses/arxiv-2603.20081-CC-BY-4.0.txt}.\par\smallskip

\noindent\textit{Editorial scope.} Counted unit: the article's proposition proposition (source0.tex lines 560--562). The article's complete original proof of it is delivered with it (source0.tex lines 564--632, one \texttt{proof} environment of 69 lines). No mathematics is written, repaired or re-derived: the statement and the proof are the article's own lines, in the article's own notation, and no retained line is substituted or re-flowed.  Retained uncounted as the source-local context the proof needs: lines 339--342; lines 544--554.  Nothing was omitted from any retained span, so the package carries no omission record.  No retained line contains a citation, so the wrapper carries no bibliography and no bibliography entry.  No retained line contains a figure environment, an image-inclusion command or drawing code, so no image is delivered and the package declares no asset dependency.  Editorial formatting only, in addition: an AMS \texttt{amsart} preamble that restates the article's own declarations and macros used by the retained lines, namely the environments \texttt{proposition} on separate counters, together with the standard packages those lines need.  The retained environments print in this document as Proposition 1 in order of appearance, whereas the article prints the same items with the numbering of its own full document.  The complete native source of the article as distributed in the arXiv e-print for this exact version is bundled unchanged as \texttt{sources/196826/source0.tex}.
\begin{equation}
\label{d:differentiable_structure}
    T_{(p_n)}\Delta = \left\{ (v_n) \in \ell^1_{\mathbb{R}} : \sum_{n=0}^{\infty} v_n = 0 \quad \text{and} \quad \left(\frac{v_n}{\sqrt{p_n}}\right) \in \ell^2_{\mathbb{R}} \right\}.
\end{equation}

\begin{equation}\label{eq:def_e_con_comp}
\bigl(\nabla^{(e)}_{(v_n)}(w_n)\bigr)_n
:=
p_n
\restr{\left(
D_{(V_n)}\!\left(\frac{W_n}{p_n}\right)
-
\sum_{k=0}^{\infty}
p_k\,D_{(V_n)}\!\left(\frac{W_k}{p_k}\right)
\right)}{(p_n)},
\end{equation}

\begin{proposition}[Exponential connection is an affine connection]
The operator \(\nabla^{(e)}\) defines an affine connection on \(T\Delta\).
\end{proposition}

\begin{proof}
Using the definition of the tangent space \(T_{(p_n)}\Delta\) in \eqref{d:differentiable_structure}, it is straightforward to check that \(\nabla^{(e)}\) is well defined, namely that
\[
\sum_{n=0}^{\infty}\bigl(\nabla^{(e)}_{(v_n)}(w_n)\bigr)_n = 0.
\]

It remains to verify that \(\nabla^{(e)}\) satisfies the axioms of an affine connection.

\textit{Linearity in the first argument.}
By the linearity of the directional derivative and \eqref{eq:def_e_con_comp}, we have
\[
\nabla^{(e)}_{f(U_n)+g(V_n)}(W_n)
=
f\nabla^{(e)}_{(U_n)}(W_n)
+
g\nabla^{(e)}_{(V_n)}(W_n)
\]
for all \(f,g\in C^{\infty}(\Delta)\) and all \((U_n),(V_n),(W_n)\in \mathfrak{X}(\Delta)\).

\textit{Linearity in the second argument.}
This follows directly from the definition of \(\nabla^{(e)}\) in \eqref{eq:def_e_con_comp}.

\textit{Leibniz rule.}
Let \(f\in C^\infty(\Delta)\). By \eqref{eq:def_e_con_comp}, together with the Leibniz rule for directional derivatives, we compute
\begin{align*}
\bigl(\nabla^{(e)}_{(V_n)}(f(W_n))\bigr)_n
&=
p_n\restr{\left(
D_{(V_n)}(f)\frac{W_n}{p_n}
+
f\,D_{(V_n)}\!\left(\frac{W_n}{p_n}\right)\right)}{(p_n)}\\
&\quad
- p_n \restr{\left(
\sum_{k=0}^{\infty}p_k\,D_{(V_n)}(f)\frac{W_k}{p_k}
+
\sum_{k=0}^{\infty}p_k\,f\,D_{(V_n)}\!\left(\frac{W_k}{p_k}\right)
\right)}{(p_n)}.
\end{align*}
Since \(D_{(V_n)}(f)\) does not depend on the summation index \(k\), and using the identity for tangent vectors from \eqref{d:differentiable_structure}, we obtain
\begin{align*}
\bigl(\nabla^{(e)}_{(V_n)}(f(W_n))\bigr)_n
&=
p_n\restr{\left(
D_{(V_n)}(f)\frac{W_n}{p_n}
+
f\,D_{(V_n)}\!\left(\frac{W_n}{p_n}\right)\right)}{(p_n)}\\
&\quad
- p_n\restr{\left(
D_{(V_n)}(f)\sum_{k=0}^{\infty}W_k
+
f\sum_{k=0}^{\infty}p_k\,D_{(V_n)}\!\left(\frac{W_k}{p_k}\right)
\right)}{(p_n)} \\
&=
\bigl(D_{(V_n)}(f)\,W_n\bigr)_n
+
f\bigl(\nabla^{(e)}_{(V_n)}(W_n)\bigr)_n \quad \forall n\in \NN.
\end{align*}
Hence
\[
\nabla^{(e)}_{(V_n)}\bigl(f(W_n)\bigr)
=
D_{(V_n)}(f)\,(W_n)
+
f\,\nabla^{(e)}_{(V_n)}(W_n)
\quad \forall f\in C^{\infty}(\Delta),\ \forall (V_n),(W_n)\in \mathfrak{X}(\Delta).
\]

Thus \(\nabla^{(e)}\) satisfies all axioms of an affine connection.
\end{proof}

\end{document}
